\section{Lifting the coarse generators modulo the uniformizer}
\label{sec:coarse-lift}

\status{The actual \texttt{GeneratorModuloUniformiser} congruence is proved here, under the source's frame and numerical data with $C_{\rm La},C_k\ge32$. The integral path basis IP3--IP4 supplies the needed parts of \texttt{NestedTypeBound}. The coarse-family dimension is a separate counting question.}

\paragraph{A common integral basis for the actual segment spaces.}
List the segment letters $\mathcal W_j$ increasingly as $x_0,\ldots,x_{N-1}$, including
the preceding frame. Write $\mathsf X_R(m,\theta,o)$ for the source's
$\mathrm{XX}'$ over a domain $R$. It is the span of separated plane/normal
products $AG$ satisfying the prefix gaps, intersected with the prescribed
normal-degree support bound. Normal variables are absent from the
preceding frame.

Apply IP3--IP4 separately to the plane and normal variables. The paired
polynomials $\Psi_{\nu,I'}\Psi_{\kappa,I''}$ form an integral basis;
the allowed subset is exactly
\[
\begin{array}{ll}
 b''_i=0,&x_i\text{ in the preceding frame},\\
 \kappa^{i+1}_1+m_{x_i}\le\nu^i_r,\quad b''_i\le o_{x_i},
       &x_i\text{ late and even},\\
 \kappa^i_1+m_{x_i}\le\nu^{i+1}_r,\quad b'_i\le o_{x_i},
       &x_i\text{ late and odd},\\
 \theta_{x_i}\le\nu^i_r,&x_i\text{ late}.
\end{array}                                                    \tag{CG.1}
\]
Indeed, expand $A,G$ using their respective lower and upper prefix
bounds. Their tensor products satisfy the gap inequalities. Projection
onto the normal-degree support condition keeps or kills each whole
paired basis vector, since its plane and normal letter degrees are fixed.
Conversely an allowed pair satisfies the defining bounds by choosing
$M_{x_i}=\kappa^{i+1}_1$ at an even letter and $\kappa^i_1$ at an odd one.
This proves the claimed basis identification.

The same integral polynomials give the basis over $O$ and its residue
field $\mathbb f$. Thus the actual $\mathsf X_O$ is saturated, reduction
maps it onto $\mathsf X_{\mathbb f}$, and projection to a total letter
multidegree preserves it. This proves the first, third and sixth
\texttt{NestedTypeBound} clauses used below and by the segment-product
argument. No integral isotypic splitting is required. \textsf{[CG1]}

\paragraph{The coarse scale and its margin.}
Put $k=r+K$, $z_j=f_{j,r-1}$ and $d_j(x)=z_j-x$. For a late letter $x$,
let $y=y_j(x)$ be its oldest of $r$ even predecessors, and set
$B_x=\deg_{y\bmod n'}/p^{d_j(y)}$, $s_x=\lfloor B_x/(k^2Y)\rfloor$.
Let $D_j$ be the product of the corresponding plane determinants to
powers $s_x$, and $\Delta=\prod_j\phi_j(D_j)$, where
$\phi_j(X_{x,c})=Z_{x\bmod n',c}^{p^{d_j(x)}}$.
Its degree is $E$, and the coarse target degree is $a'=d-E$, where
$d_t=\deg_t-E'_t$ on $TT$ and is zero elsewhere.
The source chooses
\[
 T_j=\frac{a}{8Yk^2N_s p^{d_j(f_{j-1,0})}},\qquad
 P_j=p^{R_j}\le T_j<pP_j,\qquad m_0=kN_s+1.           \tag{CG.2}
\]
The proved degree and scale estimates give
\[
 a'_t\ge0.9\deg_t\ (t\in TT),\quad E\le d,\quad P_j>m_*,\quad
 8Yk^2N_s p^{d_j(y)}P_j\le a\quad(y\in\mathcal W_j),  \tag{CG.3}
\]
where here $m_*=\texttt{mStar}$. \textsf{[CG3,RS2--RS3]}

Put $\mathcal L=C_{\rm La}N_s^2\log a$. For a coarse segment of letter degree $b$, put
$q_y=a'_{y\bmod n'}/(p^{d_j(y)}P_j)$. Its threshold and order bound are
\[
\begin{aligned}
 \theta^c_x&=\frac23
 \begin{cases}
  \lfloor q_y\rfloor-(k+r+1)-m_0,&y\text{ in the preceding frame},\\
  b(y),&y\text{ intermediate},
 \end{cases}\\
 o^c_x&=\frac{\mathcal L a}{pkp^{d_j(x)}P_j}
       +\mathbf1_{x\text{ even}}\frac{4a}{p^{d_j(y)}P_j}.
\end{aligned}                                                    \tag{CG.4}
\]
These define the coarse $\mathsf X(m_0,\theta^c,o^c)$; the actual
$M_j(b)$ also fixes degree $b$ and bounds current-frame plane exponents.
A coarse generator is
$u=Z^\ell\prod_j\phi_j(\operatorname{expand}_{P_j}y_j)$, with
$y_j\in M_j(b_j)$, physical degree $a'$, and label exponents
$\ell_{t,c}<p^{d_j(x)}P_j$ at the unique late representatives.

For a nonzero coarse generator, an intermediate physical row occurs
in only one segment. Its degree identity and its label bound give
$b(y)>q_y-k$. Set $c_y=\deg_{y\bmod n'}/(p^{d_j(y)}P_j)$.
Then $q_y\ge0.9c_y$ and $c_y\ge40kN_s$ by (CG.3).
Substitution into either case of (CG.4), using
$\lfloor q_y\rfloor>q_y-1$, proves
\[
 P_j(\theta^c_x-kN_s)\ge B_x/2,
 \qquad P_j o^c_x+kP_j\le o_x.                       \tag{CG.5}
\]
For the second inequality, the frame gap gives
$d_j(y)\ge d_j(x)+L_1+1$; hence both $kP_j$ and
$4a/p^{d_j(y)}$ are at most $a/p^{L_1+d_j(x)}$.
This retains the source's even/odd coefficients $10,1$ in $o_x$.
The companion gives both floor cases explicitly. \textsf{[CG4--CG5]}

\paragraph{The local lift.}
Split a coarse label at its unique late representative:
\[
 \ell_{t,c}=\ell^0_{t,c}+p^{d_j(x)}e_{j,x,c},\qquad
 0\le\ell^0_{t,c}<p^{d_j(x)},\quad0\le e_{j,x,c}<P_j. \tag{CG.6}
\]
For a coarse paired basis product $A_0G_0$, over $\mathbb f$ set
$A=A_0^{P_j}\ell_P D_j$ and $G=G_0^{P_j}\ell_N$, where
$\ell_P\ell_N=X^{e_j}$.
The canonical path coefficients lie in the prime field. Frobenius gives
\[
 \overline{D_j\operatorname{expand}_{P_j}(A_0G_0)X^{e_j}}=AG.
\]
Expansion fixes coefficients and is applied linearly to a general
coarse polynomial; no inverse coefficient Frobenius is assumed.

Let $d_S$ count odd plane label exponents and $d_V$ even normal label
exponents. Then $d_S+d_V\le kN_sP_j$. Take
$(\nu_*,\kappa_*)=(\nu^i_r,\kappa^{i+1}_1)$ at an even $x_i$ and
$(\nu^{i+1}_r,\kappa^i_1)$ at an odd one. At the respective prefixes,
\[
 v(A)\ge P_j\nu_*-d_S+s_x,\qquad
 M(G)\le P_j\kappa_*+d_V.                            \tag{CG.7}
\]
The determinant belonging to $x$ contributes $s_x$ because all its
rows precede $x$; every other determinant has nonnegative order.
The coarse gap $\nu_*-\kappa_*\ge kN_s+1$ and $P_j>m_*$
give the fine gap $m_*+s_x$. The always-required minimum follows from
the first inequality in (CG.5). The normal degree is at most
$P_jo^c_x+kP_j\le o_x$. Supports are unchanged. Thus $AG$ lies in the
fine $\mathsf X_{\mathbb f}$, and CG1 lifts it integrally.
Linearity handles an arbitrary coarse $y_j$; projecting its lift to
the letter degree of $D_j\operatorname{expand}_{P_j}(y_j)X^{e_j}$
preserves both membership and congruence. \textsf{[CG2,CG6--CG7]}

\paragraph{One allowed global generator.}
Call these homogeneous lifts $x_j$. The digit identity gives
\[
 \Delta u=Z^{\ell^0}\prod_j\phi_j
       (D_j\operatorname{expand}_{P_j}(y_j)X^{e_j}),
 \qquad g=Z^{\ell^0}\prod_j\phi_j(x_j).               \tag{CG.8}
\]
The segment replacements preserve all physical degrees, so $g$ has
degree $E+a'=d$, and its label and factors are allowed. Consequently
$g\in\mathrm{Gen}$ and $\overline{\Delta u}=\overline g$.
For $u=0$ take $g=0$. This closes the coarse-generator congruence
required in SD5 and SD8. \textsf{[CG8--CG9]}
